\chapter{Resultados e discussões}\label{cap:resultados}

Este capítulo apresenta a materialização do método em uma ferramenta de linha de comando e sua integração a um \textit{software} de modelagem tridimensional, bem como as limitações observadas durante os testes.

\section{A ferramenta gitbin}\label{sec:gitbin}

Em um primeiro momento, para a experimentação e a validação do método, foi utilizado um conjunto de \textit{scripts} em Python, entre eles o \texttt{glb\_parser.py}, responsável por ler e separar o arquivo \gls{glb}, e o \texttt{modificar\_glb.py}, que gerava versões com alterações controladas do modelo. Esses programas precisavam ser chamados manualmente e em uma sequência específica, de modo que era necessário conhecer a ordem e os argumentos de cada um. Antes mesmo dessa execução, também era preciso preparar o ambiente por meio de comandos como \texttt{mkfs.btrfs} e \texttt{mount -o loop}, utilizados para criar e montar a imagem do sistema de arquivos, além de criar os subvolumes e os \textit{snapshots}. Como o fluxo envolvia dezenas de etapas encadeadas, um erro em qualquer ponto do processo exigia a reconfiguração completa do ambiente.

Além disso, os caminhos utilizados eram fixos e escritos diretamente nos \textit{scripts} e nos comandos de terminal, como \texttt{/mnt/experimento} e \texttt{\textasciitilde/tcc/resultados}, o que atendia ao ambiente controlado das máquinas virtuais, mas limitava o uso do método a esse cenário. As operações também seguiam um único sentido, o que tornava trabalhoso voltar ou editar versões passadas.

Diante desses problemas, enfrentados durante a definição do método, desenvolveu-se o gitbin, uma ferramenta de linha de comando instalável via \texttt{pip} que empacota o método de versionamento diferencial em uma interface programática, eliminando a necessidade de interação direta com os comandos do \gls{btrfs} e com \textit{scripts} separados. A ferramenta está disponível sob licença aberta\footnote{\url{https://github.com/FelipeFadel/gitbin}}. Seguindo a mesma lógica de uso do Git, o gitbin implementa cinco comandos, descritos no Quadro~\ref{tab:comandos-gitbin}, e cada um deles abstrai internamente as etapas do método, como a leitura do \gls{glb}, a extração das \textit{bufferViews}, o cálculo dos \textit{hashes} \gls{sha256}, a detecção diferencial e a criação dos \textit{snapshots} imutáveis.

\begin{table}[htb]
  \centering
  \caption{\label{tab:comandos-gitbin}Comandos do gitbin e etapas do método}
  \begin{tabular}{p{2.8cm}p{6.5cm}p{4.2cm}}
    \toprule
    \textbf{Comando} & \textbf{Função} & \textbf{Etapa do método} \\
    \midrule
    \texttt{gitbin init}     & Inicializa o repositório com uma imagem \gls{btrfs} configurável & Camada física \\
    \texttt{gitbin save}     & Registra uma nova versão & Etapas 1 a 5 \\
    \texttt{gitbin status}   & Verifica as alterações pendentes & Etapas 1 a 4 \\
    \texttt{gitbin log}      & Lista o histórico de versões & \textit{Snapshots} \\
    \texttt{gitbin checkout} & Restaura uma versão anterior & \textit{Snapshots} \\
    \bottomrule
  \end{tabular}
  \fonte{Autoria própria (2026)}
\end{table}

\section{Integração com o Blender}\label{sec:addon}

Mesmo com o gitbin, o fluxo de trabalho de um artista tridimensional ainda exigia alternar entre o editor e o terminal: era necessário exportar manualmente o modelo para o formato \gls{glb}, abrir o terminal, navegar até o diretório do repositório e só então executar o comando de registro da versão. Essa troca constante de contexto representa um atrito no processo de modelagem e aumenta a chance de que versões deixem de ser registradas. Diante disso, foi desenvolvido um \textit{addon} para o Blender, [versão testada], que integra os comandos do gitbin diretamente ao editor tridimensional.

O \textit{addon} adiciona um painel lateral na \textit{viewport} 3D, apresentado na Figura~\ref{fig:painel-addon}, com os botões \textit{Save Version}, \textit{Status}, \textit{Log} e \textit{Checkout}, além de um campo para especificar a versão a ser restaurada. [O \textit{addon} aciona o gitbin como processo externo / importa o gitbin como biblioteca, e o repositório fica localizado em ...].

% \begin{figure}[htb]
%   \centering
%   \caption{\label{fig:painel-addon}Painel do \textit{addon} na \textit{viewport} 3D do Blender}
%   \includegraphics[width=0.5\textwidth]{figuras/painel-addon.png}
%   \fonte{Autoria própria (2026)}
% \end{figure}

Ao clicar em \textit{Save Version}, o \textit{addon} aciona automaticamente o exportador \gls{glb} do Blender e, em seguida, o comando \texttt{gitbin save}, materializando em um único clique o fluxo completo de exportação e versionamento. O \textit{Checkout}, por sua vez, restaura a versão informada e reimporta o arquivo resultante na cena, permitindo que o artista visualize imediatamente o estado anterior do modelo. [O \textit{addon} avisa / não avisa sobre alterações não salvas antes de restaurar uma versão.] Dessa forma, o método passou a ser utilizável em um fluxo real de modelagem, sem que o artista precise sair do editor.

Durante os testes dessa integração, entretanto, foi identificada uma limitação relevante no comportamento do exportador do Blender, descrita na seção seguinte.

\section{Limitação do exportador GLB do Blender}\label{sec:limitacao-exportador}
% Reexportação sem edição: 108 para 86 bufferViews

\section{Discussão}\label{sec:discussao}
% O que os resultados mostram, ameaças à validade